\section{Finite representatives for the original ordinary-form levels}
\label{sec:finite-levels-v42}

The ordinary projector is constructed on functions on a definite unitary
adelic quotient.  Before applying any argument about finite coefficient
rings, one needs a finite set of adelic representatives for the exact level
of those functions.  This chapter supplies that input.  It proves the full
original \texttt{FuCmpMem.finU} for arbitrary integral coefficient
representation and every permitted compact-open auxiliary level.

The standing source pin is
\texttt{7cbf7a0c32fca1badf05931879a37c3f8a2b03b4}.
All definitions below are read at that pin.  The only non-elementary input
is the already reconstructed finite-class theorem FV9, applied through an
explicit dictionary.  No integral type, ordinary structure, automorphic
Galois representation, or new published theorem is assumed.

\subsection{The objects and endpoint}

Let \(L\) be the original CM field, \(F=L^+\), and \(\sigma\) its complex
conjugation.  Let \(\ell,s,n,V,Q,c,m,\tau _0,K\) be arbitrary original
\texttt{finU} data: \(s\) is a Hida setting, \(n\ge1\),
\(V:s.\texttt{Lev}\), \(Q:\texttt{QData}\,V\),
\(\tau _0:\prod_{u\in s.S_t}\mathrm{GL}_n(\mathcal O_{L_{w(u)}})
\to\mathrm{GL}_m(c.\mathcal O)\), and \(K_q\) is a compact open
subgroup of \(\mathrm{GL}_n(L_{w_q})\) for each \(q\in Q\).
Write \(v(w)=w\cap\mathcal O_F\), \(v_q=v(w_q)\), and
\[
 S=\{v:v\mid\ell\}\cup s.S_t\cup\{s.v_0\},\qquad
 V.S'=S\cup\{v_q:q\in Q\}.
\]
The \(v_q\)'s are distinct and outside \(S\).  Put
\[
 U_{\mathbb A}=\{g\in\mathrm{GL}_n(\mathbb A_L):
                  \sigma(g)^{\mathsf t}g=1\},\quad
 U_f=\{g\in U_{\mathbb A}:g_\infty=1\},\quad
 U_\infty=\{g\in U_{\mathbb A}:g_f=1\}.
\]
Here conjugation on the full adele ring is the original
\texttt{conjAd}; the infinite coordinates are read through the original
\texttt{infEq}.  The subgroup \(\Gamma\) is the rational matrix image
\(U(F)\), literally \texttt{ratU}.  These groups are respectively
\texttt{UA}, \texttt{finU}, and \texttt{infU} in
\texttt{BaseChangeDefiniteUnitary.lean}; the name \texttt{finU} there
is a subgroup and is distinct from the endpoint theorem.

Let \(a:U_{\mathbb A}\to G^a\) be the range-restricted projection onto
all finite places of \(L\) away from \(\ell\).  For each place \(w\),
\(\operatorname{cmp}_w\) is evaluation at \(w\) when \(w\nmid\ell\)
and the trivial homomorphism otherwise.  Let \(U_b\subset G^a\) impose
integrality at the stored \(s.S_t\) places and the stored principal
congruence condition at \(w(s.v_0)\).  On this subgroup the original
\(\operatorname{sec}_w\) is the inverse of its own-place projection
on the subgroup supported above \(v(w)\), when that projection is
bijective, and is the trivial homomorphism otherwise.  Set
\[
 U_Q=\{u\in U_b:\operatorname{cmp}_w(u)\in\mathrm{GL}_n(\mathcal O_{L_w})
      \ (v(w)\notin V.S'),\quad
      \operatorname{cmp}_{w_q}(u)=1\ (q\in Q)\},
\]
\[
 U(K)=U_Q\vee\bigvee_{q\in Q}\operatorname{sec}_{w_q}(K_q),\qquad
 G^\ell=\prod_{v\mid\ell}\mathrm{GL}_n(L_{w(v)}).
\]
All joins here are subgroup joins in the source's original \(U_b\).
With \(X=(\texttt{hset}\,Q\,\tau_0).X\), the assertion to prove is
\begin{equation}\label{eq:fnl-target}
 \exists\mathcal S\subset G^a\times G^\ell\text{ finite},\quad
 \forall g\ \exists\gamma\in\Gamma,\ x\in\mathcal S,\ u\in U(K),\
 b\in X.I.\operatorname{Iw}(0):\quad
             g=\operatorname{glob}(\gamma)\,x\,(u,b).
\end{equation}
The finite set may depend on the displayed initial datum and \(K\), but
is chosen before \(g\).  It does not require a bound on its size.

% BEGIN FNL1
\paragraph{FNL1. Finite coordinate patching in an arbitrary group.}
\label{claim:FNL1}
Suppose \(I\) is finite, \(p_i:G\to H_i\) are homomorphisms and
\(P\subseteq G\) is a subgroup.  If for each \(i\) and prescribed
\(b_i\in H_i\) one has \(y_i\in P\) with
\(p_i(y_i)=b_i\) and \(p_j(y_i)=1\) for every \(j\ne i\), then there
is \(y\in P\) with \(p_i(y)=b_i\) for all \(i\).
Choose any ordering of \(I\) and multiply the \(y_i\)'s in that order.
Every coordinate of the product has exactly one possibly nonidentity
factor, so gives the required value.  Commutation of the \(y_i\)'s is
unnecessary.  For \(I=\varnothing\), use \(y=1\).
This is the original \texttt{fu\_patch} argument.
% END FNL1

% BEGIN FNL2
\paragraph{FNL2. The identity-form adelic dictionary, including topology.}
\label{claim:FNL2}
The finite unitary group for the algebraic \(F\)-group
\[
 \mathcal U(A)=\{g\in\mathrm{GL}_n(A\otimes_F L):
                     (1\otimes\sigma)(g)^{\mathsf t}g=1\}
\]
is topologically isomorphic to the literal \(U_f\).  The isomorphism
commutes with every finite-place evaluation and sends the finite image of
each rational matrix to its finite adelic extension with infinite
component \(1\).

Indeed the algebra map
\[
 \alpha:\mathbb A_{F,f}\otimes_F L\longrightarrow\mathbb A_{L,f},
 \qquad a\otimes x\longmapsto \bigl(a_{v(w)}x\bigr)_w
\]
is the original \texttt{finAdEquiv}.  Locally it is
\(F_v\otimes_F L\simeq\prod_{w\mid v}L_w\).
To check topology, choose an \(F\)-basis of \(L\).  After inverting
one nonzero integer it is a basis of an order containing and contained
in fixed multiples of \(\mathcal O_L\).  Outside finitely many places
its integral lattice agrees with the product of the local integer
rings.  Each local map and its inverse are continuous linear maps
between finite dimensional local vector spaces.  They therefore induce
inverse continuous maps of the restricted products.  Equivalently this
is the tensor topology of FV4 in these coordinates.

For a simple tensor,
\(\alpha((1\otimes\sigma)(a\otimes x))=\alpha(a\otimes\sigma x)\)
is \(\operatorname{conjFin}(\alpha(a\otimes x))\); here conjugation
fixes the embedded \(F\)-adeles and carries the \(w\) coordinate to
the conjugate place.  Linearity gives the equality on every tensor.
Applying it entrywise identifies the two unitary equations, including
the identity Gram matrix.  Extending a finite matrix by the identity
at infinity then gives \(U_f\); projection to the finite components
is its inverse.  Both maps are continuous in the displayed coordinates.
For \(x\in L\), \(\alpha(1\otimes x)\) is its diagonal finite
adelic image.  Thus rational matrices and all local projections agree
literally.  Finally, every \(g\in U_{\mathbb A}\) has a unique
factorization \(g=g_f g_\infty=g_\infty g_f\) with the indicated
finite and infinite components: extend each component by the identity
in the other factor.  This also proves that \(U_f\) and \(U_\infty\)
commute.

This dictionary checks the specialization \(J=1\), \(E=L\),
\(F=L^+\), and \(\sigma=c_L\) of FV1--FV4 and FV9.  The CM condition,
Hermitian symmetry, and nonzero determinant are automatic from these
data.  In particular no change of rational form or rational subgroup
has entered the finite-class application.
% END FNL2

% BEGIN FNL3
\paragraph{FNL3. Finite classes with the source's infinite factor in its exact position.}
\label{claim:FNL3}
For every open subgroup \(K'\subset U_f\), there is a finite
\(\mathcal T\subset U_f\subset U_{\mathbb A}\) such that
\begin{equation}\label{eq:fnl-class}
 \forall g\in U_{\mathbb A}\ \exists\gamma\in\Gamma,\ t\in\mathcal T,
 z\in U_\infty,\ k\in K':\qquad g=\gamma t z k.
\end{equation}

Apply FV9 through FNL2 to get
\(g_f=\gamma_f t k\).  For the actual rational adelic matrix
\(\gamma\), define \(z=(\gamma t)^{-1}gk^{-1}\).
Its finite component is
\(t^{-1}\gamma_f^{-1}g_f k^{-1}=1\), so \(z\in U_\infty\).
Its definition gives \(g=\gamma t z k\) with exactly that order.
The selected \(t\)'s depend only on \(K'\), not on \(g\).
This proves the full \texttt{DefU.ClassFin.classSet\_finite} conclusion
used here.  It uses the inherited checked Borel application in FV9;
it does not infer finiteness merely from finite volume.  Compactness
of \(K'\) is not required.
% END FNL3

% BEGIN FNL4
\paragraph{FNL4. All the places at which sections are needed are actually split.}
\label{claim:FNL4}
For every \(v\in V.S'\), the extension \(L/F\) is split at \(v\).
Furthermore \(v(w(v))=v\), and
\(\ell\in w\) if and only if \(\ell\in v(w)\).

The latter assertions follow by pulling back the maximal ideal of the
same stored valuation subring along
\(\mathcal O_F\to\mathcal O_L\to\Omega\); this is also NTL15.
For the split assertion, the hypothesis \(V.\texttt{split}\) says
that the absolute decomposition group at the stored place of \(\Omega\)
is contained in \(\operatorname{Gal}(\Omega/L)\).
Its image in \(\operatorname{Gal}(L/F)\) is the finite decomposition
group at \(w(v)\), hence that group is trivial.  A place in a quadratic
Galois extension has trivial decomposition group precisely when there
are two distinct places above it, each with completion \(F_v\).

Here the asserted surjectivity onto the finite decomposition group has
no unmentioned splitting assumption.  If an automorphism stabilizes
\(w(v)\), lift it to \(\Omega\); its image of the stored valuation
and the stored valuation both lie over \(w(v)\).  An element of
\(\operatorname{Gal}(\Omega/L)\) carries the former to the latter,
and correcting the lift gives an element of the absolute decomposition
group.  The conjugacy of extensions of a place follows first for finite
normal extensions from transitivity of primes, and then for \(\Omega\)
by compactness of the compatible finite Galois groups.  For completeness,
finite transitivity follows by prime avoidance and the norm: if a prime
\(\mathfrak q\) were outside the orbit of \(\mathfrak p\), choose an
integer in \(\mathfrak q\) outside every conjugate of \(\mathfrak p\).
Its norm belongs to the common base prime and is outside
\(\mathfrak p\), a contradiction.

In particular every \(w(v)\) for \(v\mid\ell\), and every added
\(w_q\), is a split place.  The proof retains the given \(V\); no
replacement set of places is chosen.
% END FNL4

% BEGIN FNL5
\paragraph{FNL5. The actual split sections and their trivial-map fallback.}
\label{claim:FNL5}
At a split place \(w\), put \(\bar w=\sigma w\).  Conjugation gives
an isomorphism \(c_w:L_w\to L_{\bar w}\).  For
\(h\in\mathrm{GL}_n(L_w)\), form the adelic unitary matrix \(j_w(h)\)
whose \(w\) coordinate is \(h\), whose \(\bar w\) coordinate is
\(c_w(h)^{-\mathsf t}\), and whose other finite and all infinite
coordinates are \(1\).  This is a continuous homomorphism.  The two
unitary equations at the split pair force exactly this relation, so
\(j_w\) is inverse to the \(w\) projection on the source subgroup
\(\operatorname{locU}_w\).  In particular it is trivial at each
place whose underlying \(F\)-place differs from \(v(w)\).

If \(v(w)\notin S\), its away projection belongs to \(U_b\): at all
the stored \(S_t\) and auxiliary-congruence coordinates it is the
identity.  Also \(w,\bar w\nmid\ell\).  Projection at \(w\) on the
source \(\operatorname{locA}_w\subset U_b\) is bijective.  Surjectivity
is provided by \(a(j_w(h))\).  For injectivity, an element in its
kernel has all coordinates away from \(v(w)\) equal to \(1\), its
\(w\) coordinate equal to \(1\), and therefore its \(\bar w\)
coordinate equal to \(1\) by the unitary relation; the away-range
coordinates separate its elements.  Thus
\[
       \operatorname{sec}_w(h)=a(j_w(h))
          \qquad(v(w)\notin S).
\]
This is an equality with the definition of the original section:
the bijective alternative of its definition has been verified and
its trivial-map fallback cannot occur at any \(w_q\).  At \(w\mid\ell\)
we use \(j_w\) in the full adelic group, not the away-section fallback.
The opposite-place inverse transpose and the distinction between these
two types of sections are essential to these identities.
% END FNL5

% BEGIN FNL6
\paragraph{FNL6. Surjectivity onto the actual frame and annihilation of the infinite group.}
\label{claim:FNL6}
The homomorphism
\[
 \Phi:U_{\mathbb A}\longrightarrow G^a\times G^\ell,
 \qquad g\longmapsto\bigl(a(g),(g_{w(v)})_{v\mid\ell}\bigr)
\]
is surjective, restricts to the original \(\operatorname{glob}\) on
\(\Gamma\), and kills \(U_\infty\).

For a prescribed \((a_0,b)\), choose \(g\) with \(a(g)=a_0\), using
that \(G^a\) is the image, not an independently specified restricted
product.  For each \(v\mid\ell\) apply FNL5 with
\(h_v=g_{w(v)}^{-1}b_v\).  Its matrix \(j_{w(v)}(h_v)\) belongs to
\(\ker a\), because both places above \(v\) divide \(\ell\).
It changes no selected \(w(v')\) coordinate with \(v'\ne v\).
The index set is finite, and FNL1 gives \(y\in\ker a\) with
\(y_{w(v)}=h_v\) for all \(v\).  Consequently \(\Phi(gy)=(a_0,b)\).
On \(\Gamma\) the displayed expression is the definition of
\texttt{hset.glob}.  An infinite-only matrix has every finite
coordinate equal to \(1\), hence \(\Phi(U_\infty)=1\).
% END FNL6

% BEGIN FNL7
\paragraph{FNL7. The exact first Iwahori subgroup is open.}
\label{claim:FNL7}
For a nonarchimedean completion \(E=L_w\), the source subgroup
\(\mathrm{I}_1(0)\) is the inverse image of the upper unipotent
subgroup of \(\mathrm{GL}_n(k_E)\) under reduction on
\(\mathrm{GL}_n(\mathcal O_E)\).  In particular it is open.  If
\(b_v\in\mathrm{I}_{1,v}(0)\) for each \(v\mid\ell\), then
\((b_v)_v\in X.I.\operatorname{Iw}(0)\).

Indeed the three defining factors of \(\mathrm{I}_1(0)\) are the
integral upper unipotents \(N\), the diagonal matrices
\(T_c(0)\) congruent to \(1\) modulo the maximal ideal, and the
lower unipotents \(N_m(0)\) congruent to \(1\).  They all reduce
to upper unipotents.  Conversely, suppose \(g\) has such reduction.
The following upper-diagonal-lower elimination proves that
\(g=x t y\) with these three types of factors.  Write
\[
 g=\begin{pmatrix}A&b\\c&d\end{pmatrix}
 =\begin{pmatrix}1&bd^{-1}\\0&1\end{pmatrix}
  \begin{pmatrix}A-bd^{-1}c&0\\0&d\end{pmatrix}
  \begin{pmatrix}1&0\\d^{-1}c&1\end{pmatrix}.
\]
Here \(d\equiv1\) is a unit, \(c\equiv0\), and the Schur
complement reduces to an upper unipotent matrix of size \(n-1\).
Induction gives the required factors; the zero-size case is trivial.
The factorization lies in the subgroup join defining \(\mathrm I_1(0)\).
Thus the claimed equality holds, and reduction is continuous to a
finite discrete matrix group.  Since the integral general linear
group is open, so is this inverse image.

Finally choose such factorizations \(b_v=x_vt_vy_v\) in each factor.
Their finite coordinatewise products satisfy
\(x\in X.I.N\), \(t\in X.I.T_c(0)\subset X.I.T\), and
\(y\in X.I.N_m(0)\).  Hence \(b=xty\in X.I.\operatorname{Iw}(0)\)
by its original subgroup-join definition.  This reconstructs the
\texttt{fu\_Iw} use of Iwahori factorization without importing the
representation-theoretic parts of \texttt{GlBridge}.
% END FNL7

% BEGIN FNL8
\paragraph{FNL8. The source's auxiliary subgroup is open on the same adelic group.}
\label{claim:FNL8}
In \(U_{\mathbb A}\) define the original subgroup
\[
 \begin{split}
 K_0=\{k:\;& k_w\in\mathrm{GL}_n(\mathcal O_{L_w})
                          \quad\hbox{for every finite }w;\\
 &k_{w(s.v_0)}\equiv1\pmod{\mathfrak m_{w(s.v_0)}};\quad
 k_{w_q}\in K_q\ (q\in Q);\quad
 k_{w(v)}\in\mathrm I_{1,v}(0)\ (v\mid\ell)\}.
 \end{split}
\]
It is open, and \(K_0\cap U_f\) is open in \(U_f\).
The all-integral condition is open: the product
\(\widehat{\mathcal O}_L=\prod_w\mathcal O_{L_w}\) is an open
subring in the finite adele ring, and membership of a matrix and its
inverse in \(M_n(\widehat{\mathcal O}_L)\) is an open condition in
\(\mathrm{GL}_n(\mathbb A_{L,f})\).  It is exactly coordinatewise
membership in the integral general linear groups.  Finite projection
and inclusion of \(U_{\mathbb A}\) are continuous by FNL2.
The principal congruence condition is the inverse image of the
identity under continuous reduction.  Each \(K_q\) is open by the
original datum and each \(\mathrm I_{1,v}(0)\) is open by FNL7.
There are only finitely many of these additional conditions, so their
intersection with the all-integral subgroup is open.  Taking its
preimage along \(U_f\hookrightarrow U_{\mathbb A}\) proves the last
assertion.  Notice that compactness of \(K_0\) was not asserted or
needed.  The auxiliary prime and every supplied \(K_q\) are unchanged.
% END FNL8

% BEGIN FNL9
\paragraph{FNL9. Transport of the actual finite subgroup into the actual level.}
\label{claim:FNL9}
For every \(k\in K_0\), write \(\Phi(k)=(u,b)\).  Then
\(u\in U(K)\) and \(b\in X.I.\operatorname{Iw}(0)\).

The \(U_b\) conditions on \(u=a(k)\) follow from the integral and
auxiliary-congruence conditions on \(k\).  If a tested coordinate
lies over \(\ell\), the source \(\operatorname{cmp}\) is the
trivial map, whose value belongs to each tested subgroup; away from
\(\ell\) it is exactly \(k_w\).  For each \(q\), put
\(h_q=k_{w_q}\in K_q\).  FNL5 gives
\(y_q=\operatorname{sec}_{w_q}(h_q)\), with its own \(q\) coordinate
\(h_q\), all other added-place coordinates equal to \(1\), and
all coordinates over places outside \(V.S'\) equal to \(1\).
Applying FNL1 inside the intersection of the indicated kernel
subgroups and \(\bigvee_q\operatorname{sec}_{w_q}(K_q)\) gives
\(y\) with all these properties simultaneously.

Set \(u_0=uy^{-1}\).  At \(w_q\) its coordinate is
\(h_qh_q^{-1}=1\); here \(w_q\nmid\ell\) follows from
\(v_q\notin S\).  At a place over \(F\setminus V.S'\), the
coordinate of \(y\) is \(1\), so \(u_0\) remains integral.
Therefore \(u_0\in U_Q\), and
\(u=u_0y\in U_Q\vee\bigvee_q\operatorname{sec}_{w_q}(K_q)=U(K)\).
This order of factors uses no unproved commutation or normality.
The assertion about \(b\) follows from its coordinatewise
\(\mathrm I_1(0)\) membership and FNL7.
% END FNL9

% BEGIN FNL10
\paragraph{FNL10. Full original finite-level theorem.}
\label{claim:FNL10}
The assertion \eqref{eq:fnl-target} holds for every original datum
specified above.  In particular the full original
\texttt{UpperTheory.V.Hi.finU} is reconstructed.

Apply FNL3 to the open subgroup \(K_0\cap U_f\), using FNL8,
and obtain a finite set \(\mathcal T\subset U_{\mathbb A}\).
Set \(\mathcal S=\Phi(\mathcal T)\).  Given any element of the
actual frame \(G^a\times G^\ell\), lift it to \(g\in U_{\mathbb A}\)
by FNL6 and write \(g=\gamma t z k\) by FNL3.  Applying \(\Phi\)
and using FNL6 and FNL9 gives
\[
 \Phi(g)=\operatorname{glob}(\gamma)\,\Phi(t)\,\Phi(z)\,\Phi(k)
         =\operatorname{glob}(\gamma)\,\Phi(t)\,(u,b),
 \quad u\in U(K),\ b\in X.I.\operatorname{Iw}(0).
\]
This is exactly the defining statement of \(X.\texttt{FinLevel}\).
The coefficient representation \(\tau _0\), its dimension \(m\),
and the coefficient ring have not been restricted at any step.
The given rank, auxiliary prime, actual section maps, Iwahori factors,
and order of multiplication have all been retained.  No condition
\(\ell>n\) is introduced; in particular rank four at \(\ell=3\)
remains in the original scope.
% END FNL10

\subsection{The remaining original Hecke-setting axioms}

The finite-level theorem also permits completion of the original
\texttt{HSet.Ax} package.  Its group-theoretic fields must refer to
the same sections as FNL5.  The arithmetic ingredients are inherited
CIP4 (the exact chosen local Artin map) and CAP6 (degree-one
Chebotarev with a prescribed Frobenius coset).  This removes an
input-producing task; it is distinct from the full ordinary/control
or Hida construction.

% BEGIN FNL11
\paragraph{FNL11. Coordinate equality, opposite places, and commuting supports.}
\label{claim:FNL11}
The family \(\operatorname{cmp}_w\) separates points of \(G^a\).
If \(v(w')=v(w)\), then
\[
 \operatorname{cmp}_w(u)=1\ \Longrightarrow\
 \operatorname{cmp}_{w'}(u)=1,\qquad
 \operatorname{cmp}_w(u)\in\mathrm{GL}_n(\mathcal O_{L_w})
 \ \Longrightarrow\
 \operatorname{cmp}_{w'}(u)\in\mathrm{GL}_n(\mathcal O_{L_{w'}}).
\]
The first assertion is the definition of the range inside a product.
For the latter assertions lift \(u\) to \(U_{\mathbb A}\).
Above a nonsplit place there is only one coordinate.  Above a split
place the other coordinate is the field conjugate inverse transpose,
which preserves the identity and integral invertible matrices.
At places above \(\ell\), both \(\operatorname{cmp}\)'s are the
trivial map, so the assertions still hold.

If \(u,x\in U_b\), \(\operatorname{cmp}_w(u)=1\), and \(x\) is
supported over \(v(w)\), then \(ux=xu\).  On the two possible
coordinates above \(v(w)\), the preceding assertion makes \(u=1\);
on every other coordinate \(x=1\).  Thus the products agree on
every coordinate.  This proves the equality in the literal subgroup,
not only in an abstract local quotient.  Two sections supported over
distinct \(F\)-places therefore commute, even before their inputs
are restricted to compact subgroups.
% END FNL11

% BEGIN FNL12
\paragraph{FNL12. Coefficient triviality and the three section commutation fields.}
\label{claim:FNL12}
Put
\(W'=\{w:v(w)\notin V.S',\ w\text{ split}\}\), the original
\texttt{HSet.Wp}.  For every added place \(w_q\) and every
\(w\in W'\), the original coefficient representation
\(\tau=\tau_0\circ\operatorname{toKT}\) is trivial on the image of
the corresponding section.  Each such section has trivial coordinates
at every place in \(s.S_t\), since its supporting \(F\)-place is
outside \(S\).  Hence its image under \(\operatorname{toKT}\) is
\(1\), and its image under \(\tau_0\) is \(1\).
This supplies both \texttt{tau\_Q} and
\texttt{coefficientRep\_trivial\_on\_unramified}.

An element of \(U_Q\) has trivial \(w_q\) coordinate, so FNL11
makes it commute with \(\operatorname{sec}_{w_q}(h)\) for every
\(h\).  Distinct added places have different underlying places by
\(Q.\texttt{inj}\), so their sections commute by FNL11.  A place
in \(W'\) lies outside \(V.S'\), while each \(v_q\) lies inside
it, so those sections commute too.  These prove exactly
\texttt{Q\_UQ}, \texttt{Q\_Q}, and
\texttt{sectionsCommuteAtSplit}.
% END FNL12

% BEGIN FNL13
\paragraph{FNL13. Integral unramified sections and the required commuting remainder.}
\label{claim:FNL13}
For every \(w\in W'\),
\[
 \operatorname{sec}_w(\mathrm{GL}_n(\mathcal O_{L_w}))\subset U_Q.
\]
At \(w\) this follows from the original inverse-section identity
FNL5.  At the conjugate place it follows from FNL11; away from the
pair the section is \(1\).  Each added-place coordinate is \(1\)
because its underlying place is in \(V.S'\).  These are all defining
conditions of \(U_Q\), proving the inclusion.

For \(u\in U_Q\) let \(h=\operatorname{cmp}_w(u)\), which is
integral, and put \(u'=\operatorname{sec}_w(h)^{-1}u\).
The inclusion just proved implies \(u'\in U_Q\), and FNL5 gives
\(\operatorname{cmp}_w(u')=h^{-1}h=1\).
Thus
\[
 u=\operatorname{sec}_w(h)u',\qquad
 [u',\operatorname{sec}_w(g)]=1\quad\hbox{for every }g,
\]
where the second equality is FNL11.  This is the original
\texttt{W\_dec} order, and the preceding inclusion is
\texttt{W\_le}.
% END FNL13

% BEGIN FNL14
\paragraph{FNL14. The same-place alternative preserves the actual sections and maximal compact groups.}
\label{claim:FNL14}
Let \(w,w'\in W'\).  If their underlying \(F\)-places differ,
all their section images commute.  If their underlying places agree,
there is a group isomorphism
\[
 e:\mathrm{GL}_n(L_{w'})\longrightarrow\mathrm{GL}_n(L_w)
\]
carrying the integral general linear group onto the integral general
linear group and satisfying
\(\operatorname{sec}_{w'}(g)=\operatorname{sec}_w(e(g))\) for every
\(g\).  For equal places take the identity.  For the two conjugate
places take the field-conjugate inverse-transpose isomorphism.
FNL5 identifies both sides coordinatewise, and FNL11 verifies the
integral equivalence in both directions.  This proves the exact
\texttt{Wm\_mul\_Wm\_comm} alternative; it never assumes that two
opposite-place sections commute.
% END FNL14

% BEGIN FNL15
\paragraph{FNL15. Density of inverse Frobenius on the original quotient.}
\label{claim:FNL15}
Let \(G_0=\ker\eta\subset G_{F,S'}\) be the original quotient group
in \(V.\texttt{frame}\), and
\(p=V.\operatorname{pr}:G_L\to G_0\) the original map.  It is a
continuous surjection.  Moreover
\begin{equation}\label{eq:fnl-dense}
 \{p(\sigma^{-1}):w\in W',\ A\text{ lies over }w,\
       \sigma\text{ is arithmetic Frobenius at }A\}
                    \quad\text{is dense in }G_0.
\end{equation}

The original quotient map \(q:G_F\to G_{F,S'}\) is continuous and
onto.  Its sign character has
\(\eta(q(g))=1\) exactly when \(g\) fixes \(L\), by the stored
\(V.\texttt{h}\eta\) identity.  A lift under \(q\) of any element
of \(G_0\) therefore belongs to \(G_L\), proving surjectivity of
\(p\).  Continuity follows from continuity of restriction to the
same algebraic closure, the quotient topology, and the subgroup
topology.  Consequently \(g\mapsto p(g^{-1})\) is continuous and
surjective as well.

Only finitely many places of \(L\) lie over \(V.S'\): there are
finitely many base places, each with finitely many factors in the
number field extension.  Add the finitely many places ramified in
\(L/F\) to this excluded set.  CAP6, applied with \(E=L\), says
that arithmetic Frobenius elements at the remaining degree-one
places meet every open coset of \(G_L\).  Such cosets form a basis
of its profinite topology, so those Frobenius elements are dense.
For one of these places \(w\), absolute degree one means
\(k_w=\mathbb F_p\).  Hence also \(k_{v(w)}=\mathbb F_p\) and the
relative residue degree is one.  The relative ramification index is
one by our exclusion.  The quadratic Galois degree formula
\(2=e f g\) then gives two primes above \(v(w)\), so \(w\) is split.
It is therefore in \(W'\).

The image of a dense subset under a continuous surjection is dense:
every nonempty open set has a nonempty open preimage.  Apply this
to \(g\mapsto p(g^{-1})\) and the Frobenius set just obtained.
This proves \eqref{eq:fnl-dense} with the source's exact inverse
and its arbitrary valuation-over-place witnesses.  CAP6 supplies
both those witnesses and the actual arithmetic Frobenius action;
no density assertion about an unrelated quotient is substituted.
% END FNL15

% BEGIN FNL16
\paragraph{FNL16. The complete original \texttt{HSet.Ax} on the constructed setting.}
\label{claim:FNL16}
For every original datum in FNL10, the setting
\(H=\texttt{hset}\,Q\,\tau_0\) satisfies all eighteen fields of
\(H.\texttt{Ax}\).

Here is the complete assembly, retaining the original field names.
\begin{center}\small
\begin{tabular}{p{0.43\linewidth}p{0.49\linewidth}}
\texttt{odd}, \texttt{pos} & The stored \(s.\texttt{odd}\) and \(n\ge1\).\\
\texttt{Sl\_over}, \texttt{Sl\_inj} & FNL4, by contraction of places.\\
\texttt{art\_spec} & CIP4 on the same stored valuation and \(\operatorname{artOf}\).\\
\texttt{Q\_away}, \texttt{Wp\_away} & The defining place exclusions and FNL4.\\
\texttt{fin} & FNL10 for every permitted \(K\).\\
\texttt{tau\_Q}, \texttt{coefficientRep\_trivial\_on\_unramified} & FNL12.\\
\texttt{Q\_UQ}, \texttt{Q\_Q}, \texttt{sectionsCommuteAtSplit} & FNL12.\\
\texttt{W\_le}, \texttt{W\_dec} & FNL13.\\
\texttt{Wm\_mul\_Wm\_comm} & FNL14.\\
\texttt{prj\_cont}, \texttt{dense} & FNL15.
\end{tabular}
\end{center}
In particular \(w(v)=w(v')\) implies \(v=v'\) by contraction,
which verifies \texttt{Sl\_inj} without a choice of new embeddings.
The local Artin field uses CIP4's validity of the original Hilbert-choice
map, so no further \texttt{LiteratureRankNHypothesis} premise is needed
in this human proof beyond its already proved class-field-theory input.
This is the full original \texttt{axU} conclusion and, in particular,
the conclusion of \texttt{CLocConj.axU\_s}.  The latter source variant
carries a larger literature package because of its Artin-map route;
CIP4 is a proved replacement for that use.
% END FNL16

\subsection{What this supplies, and what still has to be constructed}

Together with NTL17 and OTW12, FNL10 supplies finite level and the actual
\texttt{Frame.Free} stabilizer condition needed at the entrance to
ordinary theory.  FNL16 supplies the entire original
\texttt{HSet.Ax}, including inverse-Frobenius density, so this package
is no longer an unconstructed antecedent.  The argument does not by
itself supply ordinary Hecke commutation, classical control,
arithmetic-point separation, the integral type, or same-ideal all-level
automorphy.  The outer \texttt{COUNT\_II\_SUPPORT} obligation remains
open until those producers have been proved.  The finite-class,
local-reciprocity and Chebotarev results are used as completed
mathematical inputs with their actual objects and maps checked.
